\section{Baselines, External Benchmarks, Real Robot, and Timing}
\label{app:evaluation}

\subsection{Baselines}
\label{app:baselines}
\paragraph{Common conditions.} Table~\ref{tab:baselines} summarizes the adaptations.
Each baseline keeps its official architecture, inputs, action chunk, and
initialization. ACT and Diffusion Policy train from scratch, $\pi_{0.5}$ is
fine-tuned from \texttt{pi05\_base}, and the other methods start from their released
weights. All of them train on the same 11{,}200 demonstrations, whose commands we
map to each action space through the same Cartesian controller. \bench renders a
$224\times224$ wrist camera in addition to the agent view, and methods whose recipes
use a wrist view receive it. \ours and the kinematic baselines use the agent view
only. All methods run under the shared executor and the real-time protocol of
Appendix~\ref{app:sweeps}, and Table~\ref{tab:supp_protocols} adds their original
executors.

\paragraph{Training.} Each baseline follows its official recipe for the published
number of updates and is trained once, with seed 0. We changed only the data
loading and the action space. For DreamVLA, we ran both its original 40-epoch
schedule and a 20-epoch cosine schedule and report the one that was better on
development scenes. The configuration of \ours, in contrast, was selected by the
ablations of Appendix~\ref{app:design}, a search that the baselines do not receive,
and its variants use three training seeds.

\paragraph{Inputs.} PUMA \cite{domino} receives its historical optical flow, computed
on the rendered frames. DreamVLA receives seven causal RGB frames, the end-effector
position and orientation, the finger positions, and a fixed task text, and it samples
with ten DDIM steps and classifier-free guidance of 1.5. OpenWAM-$\alpha$ receives
the current RGB views, joint values, gripper opening, and task text, and it generates
video and actions jointly with video decoding disabled. DynamicWAM receives its flow
history and kinematic descriptors.

\paragraph{SE(3)-augmented baselines.} $\pi_{0.5}$+SE(3) and DynamicWAM+SE(3)
receive the eight target-state tokens of Appendix~\ref{app:tokens} from the same
pose source, with the same validity flags and missing-state embedding. A two-layer
MLP maps the tokens to the model's width. For $\pi_{0.5}$, we append them to the
prefix after the image tokens, where the action expert attends to them, and for
DynamicWAM they join its kinematic descriptors. The output layer of the MLP is
zero-initialized, so training starts from the behavior of the base model, which
also provides the recipe.

\paragraph{Kinematic baselines.} \emph{\ours, CV input} shares the scene tokens and
the action stream of \ours but has no forecast losses. Instead, the action stream
receives the constant-velocity extrapolation of the target pose at the execution
times of its chunk as input tokens, and the model is trained on the same data with
the same recipe. Appendix~\ref{app:design_predict} adds a variant that receives these
tokens and keeps the forecast losses. The \emph{Kalman interceptor} has no learned
components. A constant-acceleration Kalman filter \cite{barshalom2001} runs on the
target pose, with its process noise tuned on development scenes, and predicts the
target's trajectory. The interceptor chooses the earliest point of this trajectory
that the arm can reach under its velocity limits after its measured latency,
approaches it along the expert's pre-grasp offset, and closes, releases, or aligns
at the predicted time. It replans at every query and uses the expert's
task-specific grasp, placement, and entry frames.

% Protocol table: Dynamic-14 baseline adaptations from BASE1 and the v3 endpoint records.
% PUMA uses the resolved 40k training configuration and parameter summary in
% BASE1/results/puma_dynamic14_20261101/h100/training-final.zip.
\begin{table}[!htbp]
\centering\footnotesize
\caption{\tablelang{\textbf{Baseline training.} The training budget gives optimizer updates and global batch size.
$H$ is the number of commands returned per query. The parameter counts include trainable parameters only.
DOM ckpt. denotes the released DynamicVLA checkpoint trained on DOM.}{\textbf{基线训练。}训练量以优化器更新次数和全局 batch 大小表示。
$H$ 为每次查询返回的指令数。参数量只计入可训练参数。DOM 检查点指 DynamicVLA 发布的、在 DOM 上训练的检查点。}}
\label{tab:baselines}
\setlength{\tabcolsep}{2.6pt}
\begin{tabular}{@{}llccc@{}}
\toprule
\tablelang{Method}{方法} & \tablelang{Initialization}{初始化} & \tablelang{Updates$\times$batch}{更新数$\times$batch} & $H$ & \tablelang{Params}{参数量} \\
\midrule
ACT \cite{act} & \tablelang{scratch}{从头训练} & $50{,}000\times32$ & 100 & 83.9M \\
DP \cite{diffusionpolicy} & \tablelang{scratch}{从头训练} & $50{,}000\times32$ & 15 & 288.7M \\
$\pi_{0.5}$ \cite{pi05} & \texttt{pi05\_base} & $30{,}000\times256$ & 10 & 3,353.4M \\
$\pi_{0.5}$+SE(3) & \texttt{pi05\_base} & $30{,}000\times256$ & 10 & 3,357.7M \\
DynamicVLA \cite{dynamicvla} & \tablelang{DOM ckpt.}{DOM 检查点} & $50{,}000\times44$ & 20 & 99.1M \\
PUMA \cite{domino} & \tablelang{official}{官方} & $40{,}000\times64$ & 16 & 4,516.5M \\
OpenWAM-$\alpha$ \cite{openwam_alpha} & \tablelang{official}{官方} & $34{,}490\times192$ & 32 & 6,021.2M \\
DynamicWAM \cite{dynamicwam} & \tablelang{official}{官方} & $40{,}000\times128$ & 16 & 988.8M \\
DynamicWAM+SE(3) & \tablelang{official}{官方} & $40{,}000\times128$ & 16 & 989.4M \\
\midrule
\ours & \tablelang{scratch}{从头训练} & $\cfgUpdates\times512$ & \cfgChunk & \cfgTrainableParams \\
\bottomrule
\end{tabular}
\end{table}


\subsection{DOMINO and LIBERO}
\label{app:external}
DOMINO \cite{domino} contains 35 tasks of two interaction types, interception and
tracking, at three dynamics levels: constant velocity, polynomial trajectories, and
stochastic segments with abrupt changes. We use its \fillin{N} single-arm tasks on
the Franka embodiment in the clean setting, at a speed coefficient of \fillin{0.1}
and at all three levels, with the official demonstrations, the benchmark's own
evaluator, and 100 closed-loop episodes per task and training seed. We report SR and
DOMINO's manipulation score (MS), which scales route completion by penalties for
leaving the workspace or the view and for collisions. On LIBERO \cite{libero}, we use
the Spatial, Object, Goal, and Long suites with ten tasks each, 50 evaluation
episodes per task, three training seeds for our variants and one for baselines, and
the released demonstrations and initial states. LIBERO checks that the role-aware
forecast does not cost performance on static tasks. On both benchmarks, \ours trains
from the same generic initialization as on \bench, with domain-specific
normalization. As on \bench, it reads the pose of the object that the instruction
acts on first from the simulator, and queries without a valid pose
(\fillin{XX}\% on LIBERO and \fillin{XX}\% on DOMINO) use the missing-state
embedding of Appendix~\ref{app:tokens}. The pose is privileged information that the
baselines do not receive, so Table~\ref{tab:supp_external} also reports \ours with
tracker-calibrated pose noise. \ours reaches \fillin{XX.X}\% on DOMINO against
\fillin{XX.X}\% for \fillin{the strongest DOMINO baseline}, and on LIBERO it is
within \fillin{X.X} points of $\pi_{0.5}$.

% Planned results: A8 external comparison; four methods, one training seed each, 50 episodes per task and method.
\begin{table*}[t]
\centering\small
\caption{\tablelang{\textbf{DOMINO evaluation} on 24 single-arm clean-L1 tasks (one training seed per method, 50 episodes per task and method). SR (\%) is grouped by interaction type and averaged over tasks; MS is the manipulation score. Protocol: Appendix~\ref{app:external}.}{\textbf{DOMINO 评测}（24 个单臂 clean-L1 任务；每个方法一个训练种子，每个任务与方法 50 个回合）。SR（\%）按交互类型分组并在任务上平均；MS 为操作得分。协议见附录~\ref{app:external}。}}
\label{tab:supp_external}
\setlength{\tabcolsep}{7pt}
\begin{tabular}{@{}lcccc@{}}
\toprule
\tablelang{Method}{方法} & \tablelang{Intercept}{拦截} & \tablelang{Track}{跟踪} & \tablelang{All}{全部} & MS \\
\midrule
$\pi_{0.5}$ & \resultplaceholder & \resultplaceholder & \resultplaceholder & \resultplaceholder \\
DynamicWAM & \resultplaceholder & \resultplaceholder & \resultplaceholder & \resultplaceholder \\
PUMA & \resultplaceholder & \resultplaceholder & \resultplaceholder & \resultplaceholder \\
\ours & \resultplaceholder & \resultplaceholder & \resultplaceholder & \resultplaceholder \\
\bottomrule
\end{tabular}
\end{table*}


\subsection{Real Robot}
\label{app:real}
\begin{figure}[t]
\centering
\figplaceholder[1.6in]{\linewidth}{Fig.~S4 --- real-robot setup}{%
\textbf{Left:} photo of the cell: Franka arm with gripper, fixed third-person RGB-D camera, wrist camera, and the devices that move the
targets (\fillin{belt conveyor, turntable, ramp}), with the base frame drawn. \textbf{Right:} one frame per task at the instant of the first
command, annotated with the target's motion type (linear, circular, accelerating) and $\|\vect v\|\Delta$ for the measured latency.}
\caption{Real-robot setup: the cell and the five tasks.}
\label{fig:real_setup}
\end{figure}

\paragraph{Platform.} The robot is a seven-DoF Franka arm with a parallel-jaw
gripper, commanded through a timestamped Cartesian interface at \cfgControlRate{}
(Fig.~\ref{fig:real_setup}). A fixed \fillin{RealSense D435} records RGB-D at
\fillin{$640\times480$} and \fillin{30}\,Hz, and we calibrate its extrinsics with a
checkerboard before each session. \ours uses this third-person view, and
$\pi_{0.5}$ also receives the image of a second camera on the wrist, as its recipe
requires.

\paragraph{Tasks and success.} The five tasks cover three kinds of target motion:
linear motion (\fillin{T1, T2}), circular motion (\fillin{T3, T4}), and accelerating
motion (\fillin{T5}). \fillin{For each task: the device that moves the target, the
speed range, the objects, and the success condition, e.g., both fingers hold the
object stably and lift it for at least 0.5\,s.} Touching or pushing the target, a
single-finger hold, and an action after the target has left the workspace count as
failures. We measure the speed of every moving device at the start and end of every
block, and trials whose speed deviates from the nominal value by more than 5\% are
repeated and reported as equipment failures.

\paragraph{Trials.} Each method runs \fillin{20} trials per task. All methods use
the same pre-generated trial keys, which fix the object instance, the speed, and the
start pose, sampled from the calibrated reachable region. An external fixture that
the policy does not see sets and records the start pose. Each task's trials run in
random order, so that drift of the apparatus does not align with the methods. An
operator resets the scene by hand. A trial starts at the first usable observation,
without a free observation period, and the robot may move at any time. Human
contact, safety stops, and manual completion count as failures. Infrastructure
failures before the first usable observation are repeated and reported. Test trials
use object instances that appear in no demonstration, and demonstrations,
development trials, and test trials come from separate sessions.

\paragraph{Fine-tuning.} \ours is fine-tuned from the DROID base on \fillin{N}
teleoperated demonstrations per task (Appendix~\ref{app:posttraining}). The
reference without the base starts from the generic initialization with the same
demonstrations and search, and $\pi_{0.5}$ and DynamicWAM are fine-tuned from their
released weights on the same demonstrations. The tracker uses \fillin{a CAD model or
reference views of each object} (Table~\ref{tab:supp_tracker}).

\paragraph{Analysis.} We measure the contact-time error against the feasible
contact window computed from the device's encoder, and failures without contact
remain in the SR denominator. Table~\ref{tab:supp_real_ci} reports successes per
task with 95\% binomial intervals for the pooled rates and paired differences from a
bootstrap over trial keys. The evaluation took \fillin{X} days. Most failures of
\ours are \fillin{late closures after tracking losses under the gripper}.

% Planned results: six-task real-robot study (App. D.3); per-task counts, baseline set, and numeric results to be filled (see docs/writing/v4_t5_real_report.md).
\begin{table*}[t]
\centering\footnotesize
\caption{\tablelang{\textbf{Real-robot successes} per task (successes/trials, \fillin{n} trials per task; tasks in Table~\ref{tab:real_tasks}). All: pooled success rate (\%) with its 95\% binomial interval. $\Delta$: paired difference from \ours (points) with a bootstrap interval over trial keys. $e_t$: median contact-time error (ms) on trials that make contact.}{\textbf{真机成功次数}，按任务分列（成功次数/试验数，每个任务 \fillin{n} 次；任务定义见表~\ref{tab:real_tasks}）。All：合并成功率（\%）及 95\% 二项区间。$\Delta$：相对 \ours 的配对差值（百分点）及按试验键自助法得到的区间。$e_t$：发生接触的试验中接触时刻误差的中位数（ms）。}}
\label{tab:supp_real_ci}
\setlength{\tabcolsep}{3pt}
\begin{tabular}{@{}lccccccccc@{}}
\toprule
\tablelang{Method}{方法} & T1 & T2 & T3 & T4 & T5 & T6 & All & $\Delta$ & $e_t$ \\
\midrule
\fillin{$\pi_{0.5}$} & \resultplaceholder[s/n] & \resultplaceholder[s/n] & \resultplaceholder[s/n] & \resultplaceholder[s/n] & \resultplaceholder[s/n] & \resultplaceholder[s/n] & \resultplaceholder[00.0 {[00, 00]}] & \resultplaceholder[{$-$00.0 [00, 00]}] & \resultplaceholder \\
\fillin{DynamicWAM} & \resultplaceholder[s/n] & \resultplaceholder[s/n] & \resultplaceholder[s/n] & \resultplaceholder[s/n] & \resultplaceholder[s/n] & \resultplaceholder[s/n] & \resultplaceholder[00.0 {[00, 00]}] & \resultplaceholder[{$-$00.0 [00, 00]}] & \resultplaceholder \\
\tablelang{\ours, no DROID base}{\ours，无 DROID 基座} & \resultplaceholder[s/n] & \resultplaceholder[s/n] & \resultplaceholder[s/n] & \resultplaceholder[s/n] & \resultplaceholder[s/n] & \resultplaceholder[s/n] & \resultplaceholder[00.0 {[00, 00]}] & \resultplaceholder[{$-$00.0 [00, 00]}] & \resultplaceholder \\
\ours & \resultplaceholder[s/n] & \resultplaceholder[s/n] & \resultplaceholder[s/n] & \resultplaceholder[s/n] & \resultplaceholder[s/n] & \resultplaceholder[s/n] & \resultplaceholder[00.0 {[00, 00]}] & -- & \resultplaceholder \\
\bottomrule
\end{tabular}
\end{table*}


\subsection{Timing}
\label{app:timing}
We time all methods on one \cfgGPU{} (32\,GB), the robot's deployment GPU. For each
method, timing uses thirty warm-up calls and 300 measured calls over three
observation sequences of eight frames, with CUDA synchronization at call boundaries
and three rounds with rotated method order, and each method uses the precision of
its release. The time-to-action includes encoding the image and generating the
first command, and on the robot also tracking the newest frame. For \ours, we
record the time of the forecast heads separately, although deployment skips them.
Table~\mainref{tab:main} gives the p95 values and Table~\ref{tab:baselines} the
parameter counts. On the robot, traces record capture, availability,
preprocessing, the start and end of inference, queue entry, command issue, and
execution. The action age is the execution time minus the capture time, and we
report its p50, p95, and maximum together with the fraction of missed deadlines and
the peak memory (Table~\ref{tab:supp_timing}). The latency sweep of \bench differs
from wall-clock deployment, because it inserts a prescribed delay while physics
advances.

% Transcribed from results/realtime_c/stage-timing-20261104.json: accepted compiled BF16 C stage measurement, separate eager readouts and unprobed totals. Body only.
\caption{\tablelang{\textbf{Latency breakdown} of compiled BF16 \ours on one \cfgGPU{} at batch size 1 (ms, p50 / p95). Stage rows are CUDA-event spans over 900 queries; totals are uninstrumented wall time, excluding pose tracking. The original Table~\mainref{tab:main} calibration is retained alongside this run's control. Readouts are 900 separate BF16 eager forwards on one fixed recorded World feature; they are absent from deployment. Stage percentiles are not additive. Magenta dashes mark pending robot measurements.}{\textbf{延迟分解}（BF16编译版 \ours，\cfgGPU{}，batch 为1，ms；p50 / p95）。阶段项为900次查询的CUDA事件区间，合计为无探针墙钟时间，不含位姿跟踪；同时保留表~\mainref{tab:main} 的原始标定与本轮对照。仅训练读出在同一固定录制World特征上单独执行900次BF16 eager前向，不参与部署。阶段分位数不相加；品红色破折号表示待测的真机项。}}
\label{tab:supp_timing}
\setlength{\tabcolsep}{4pt}
\begin{tabular}{@{}lcc@{}}
\toprule
\tablelang{Stage}{阶段} & p50 & p95 \\
\midrule
\tablelang{Pose tracking (per frame, robot)}{位姿跟踪（每帧，真机）} & \resultplaceholder & \resultplaceholder \\
\tablelang{Observation encoding + H2D}{观测编码与传输} & 4.83 & 5.28 \\
\tablelang{World/KV preparation (once/query)}{World/KV准备（每次查询一次）} & 3.93 & 4.22 \\
\tablelang{Action generation (\cfgEuler{} Euler steps)}{动作生成（\cfgEuler{} 个Euler步）} & 4.27 & 4.41 \\
\tablelang{Policy total (Table 1 calibration)}{策略合计（表1标定）} & 14.59 & 15.11 \\
\tablelang{Policy total (unprobed control)}{策略合计（本轮无探针对照）} & 14.66 & 15.65 \\
\tablelang{Training-only readouts (separate, eager)}{仅训练读出（单独、eager）} & 0.37 & 0.41 \\
\tablelang{Action age at execution (robot)}{执行时的动作年龄（真机）} & \resultplaceholder & \resultplaceholder \\
\tablelang{Queries over 50\,ms (robot, \%)}{超过 50\,ms 的查询（真机，\%）} & \multicolumn{2}{c}{\resultplaceholder} \\
\bottomrule
\end{tabular}

